\scanpage{173}
\mspage{8}
% archive sheet number: 172
Il n'est pas vrai en général que \deleted{\unsure{toujours}} ce torseur à opérateurs $I$ soit trivial, i.e. admette un élément invariant par $I$;
\deleted{Mais si $P$ est le groupe de Sylow, on}
\inserted{l'obstruction est dans un $H^1(I,\pi_1(Y_{\bar K},\bar a))$}
($\unsure{H^1|H}$ non commutatif). Mais considérons le \deleted{\unsure{$p$}} torseur $Q^{\unsure{(p')|(p)}}$ associé de groupe \deleted{$\unsure{R|K}$} structural (à opérateurs $I$) $\pi_1^{(p')}(Y_{\bar K},\bar a)$, et la restriction des $I$-opérations au s-groupe $P\subset I$ ($p$-sous-groupe de Sylow). Je dis qu'il existe un élément de $Q^{(p')}$ fixe sous $P$, i.e. que la classe $c^{(p')}|P\in H^1(P,\pi_1^{(p')}(Y_{\bar K},\bar a))$ est nulle : Ceci provient en effet du fait que $P$ est un $p$-groupe et \inserted{$G:$} $\pi_1^{(p')}(Y_{\bar K},\bar a)$ est premier à $p$ : à prouver que \deleted{\unsure{toute}} toute section multiplicative $\varphi$ du produit $\frac12$ direct $G.I\to I$ est conjuguée de la section \unsure{standard|standart} par un élément de $G$. Or par le th de \inserted{conjugaison} Sylow, il existe un élément $g.i\in G.I$ tel que $\operatorname{int}(g.i)I=\varphi(I)$, i.e. $\operatorname{int}(g)I\subset\varphi(I)$, d'où la conclusion.
